\subsection{Biểu diễn thời gian bằng Temporal Encoding}
\label{subsec:temporal_encoding}

Bên cạnh cấu trúc không gian của mạng lưới giao thông, tín hiệu giao thông còn thể hiện tính
chu kỳ mạnh theo thời gian, đặc biệt là theo \textit{thời điểm trong ngày} (giờ cao điểm sáng/chiều)
và \textit{ngày trong tuần}. Do đó, để mô hình Transformer có thể phân biệt được các mẫu hình
giao thông theo ngữ cảnh thời gian và khai thác được tính tuần hoàn, nhóm thiết kế một cơ chế
\textit{Temporal Encoding} để mã hóa thông tin thời gian và bổ sung vào embedding đầu vào.

\subsubsection{Mã hóa vị trí tương đối theo bước thời gian}
\label{subsubsec:relative_pos}

Đối với mỗi mẫu huấn luyện, dữ liệu đầu vào được cắt theo một cửa sổ thời gian dài $L$ bước:
$\{t-L+1,\dots,t\}$. Vì Transformer không có khái niệm thứ tự nội tại giữa các token, nhóm sử dụng
\textit{relative position embedding} để mã hóa vị trí của mỗi bước thời gian trong cửa sổ này.
Cụ thể, với chỉ số vị trí tương đối $p \in \{0,1,\dots,L-1\}$, mỗi vị trí được gán một vector embedding
học được:
\begin{equation}
\mathbf{e}_p \in \mathbb{R}^{d_{\text{model}}}.
\end{equation}
Cơ chế này giúp mô hình phân biệt rõ các quan sát gần hiện tại và các quan sát xa hơn trong quá khứ,
từ đó học được động lực thay đổi của giao thông theo thời gian.

\subsubsection{Mã hóa chu kỳ theo thời điểm trong ngày và ngày trong tuần}
\label{subsubsec:cyclic_time_features}

Để khai thác tính chu kỳ của giao thông, nhóm sử dụng các đặc trưng thời gian gồm:
(i) thời điểm trong ngày đã chuẩn hóa $\text{tod}(t) \in [0,1]$, và
(ii) ngày trong tuần $\text{dow}(t) \in \{0,\dots,6\}$.
Các đặc trưng này được đưa về dạng biểu diễn lượng giác nhằm phản ánh tính tuần hoàn:
\begin{align}
\mathbf{c}_{\text{tod}}(t) &= \big[\sin(2\pi\,\text{tod}(t)),\ \cos(2\pi\,\text{tod}(t))\big],\\
\mathbf{c}_{\text{dow}}(t) &= \big[\sin(2\pi\,\text{dow}(t)/7),\ \cos(2\pi\,\text{dow}(t)/7)\big].
\end{align}
Sau đó, nhóm ghép các thành phần trên thành một vector 4 chiều và chiếu tuyến tính sang không gian
$d_{\text{model}}$:
\begin{equation}
\mathbf{g}(t) = \mathbf{W}_t \big[\mathbf{c}_{\text{tod}}(t)\ \Vert\ \mathbf{c}_{\text{dow}}(t)\big] + \mathbf{b}_t,
\quad \mathbf{W}_t \in \mathbb{R}^{d_{\text{model}} \times 4}.
\end{equation}

\subsubsection{Kết hợp Temporal Encoding vào embedding đầu vào}
\label{subsubsec:temporal_integration}

Đối với mỗi bước thời gian trong cửa sổ đầu vào, Temporal Encoding được tạo bằng cách cộng hai thành phần:
(1) embedding vị trí tương đối theo chỉ số $p$ và (2) embedding chu kỳ theo ngữ cảnh thời gian thực:
\begin{equation}
\mathbf{t}(p) = \mathbf{e}_p + \mathbf{g}(t-L+1+p),
\quad p = 0,\dots,L-1.
\end{equation}
Vector $\mathbf{t}(p)$ sau đó được cộng vào embedding của tất cả các tuyến đường tại bước thời gian $p$.
Nhờ đó, cơ chế self-attention có thể điều kiện hóa trên ngữ cảnh thời gian, giúp mô hình học được các mẫu
hình giao thông mang tính lặp lại (ví dụ cao điểm sáng/chiều) và phân biệt các trạng thái giao thông giữa
các ngày trong tuần.

Tóm lại, Temporal Encoding cung cấp cho mô hình Transformer một tín hiệu thời gian giàu ngữ cảnh,
bổ sung cho Spatial Encoding dựa trên Laplacian, tạo nền tảng cho việc học biểu diễn không gian--thời gian
trong phần kiến trúc mô hình trình bày ở mục tiếp theo.
